\chapter{Glossary and cheat sheet}

\section{Terms in plain words}

\begin{description}[leftmargin=3.4cm,style=nextline,font=\bfseries]
\item[AC coefficient] Any DCT output other than the first: how much of a particular wiggly pattern is present.
\item[Amplitude bits] The raw bits after a Huffman code that say which value inside a category it is.
\item[bpp] Bits per pixel of the compressed file.
\item[Canonical Huffman] A Huffman code whose bit patterns are assigned by a fixed rule from the code lengths, so only the lengths need be stored.
\item[Category / size] Number of bits needed for the magnitude of a value.
\item[Coefficient] One of the 64 numbers a tile is turned into by the DCT.
\item[DC coefficient] The first DCT output: 8$\times$ the tile's average brightness.
\item[DCT] Discrete Cosine Transform: re-expresses a tile as a weighted sum of 64 cosine patterns.
\item[DPCM] Sending differences to a prediction (here: previous tile's DC).
\item[EOB] End of block: ``all remaining coefficients are zero''.
\item[Entropy] The theoretical minimum average bits per symbol for a given distribution.
\item[Fused / staged] Two ways to run the same stages: per tile in a tiny buffer, or as full-image passes with big intermediates.
\item[Kraft inequality] $\sum2^{-\ell_i}\le1$ for any prefix-free code; equality means no wasted codewords.
\item[Level shift] Subtracting 128 so pixels are centred on zero.
\item[Lossy / lossless] Whether the original can be recovered exactly.
\item[MSE / PSNR] Average squared error / its logarithmic form in dB.
\item[Orthonormal] A transform that preserves lengths and angles; inverse = transpose.
\item[Parseval] Orthonormal transforms preserve the sum of squares.
\item[Prefix-free] No code is the beginning of another.
\item[Quality] The 1--100 slider that scales the quantization table.
\item[Quantization] Dividing by a step and rounding: the lossy step.
\item[Rate--distortion] The trade-off between file size and error.
\item[Run-length coding] Describing zeros by how many there are.
\item[Slurm] The job scheduler on bwUniCluster.
\item[Tile / block] An $8\times8$ group of pixels.
\item[ZRL] Zero run length: ``16 zeros'', used when a run exceeds 15.
\item[Zig-zag] Diagonal scan order from low to high frequency.
\end{description}

\section{Command cheat sheet}

\begin{lstlisting}[style=shell]
# local
make && make test                               # build, 28 groups / 32,581 checks
python3 python/make_test_images.py              # data/*.pgm
./build/jlc eval data/camera.pgm -q 50          # one round trip with metrics
python3 python/crosscheck.py                    # C++ vs NumPy, expect 0 mismatches
python3 python/rd_sweep.py                      # rate-distortion
python3 python/numerics_study.py                # numerics
./build/jlc_bench --sizes 512,1024,2048 --quality 10,50,90 --reps 7 --fused 0
# cluster (from the project directory on uc3.scc.kit.edu)
mkdir -p logs results data
sbatch slurm/01_build_test.sbatch               # dev_cpu: build + tests
bash slurm/submit_all.sh                        # benchmarks on cpu_il and cpu
sbatch slurm/03_profile.sbatch                  # gprof
squeue --me                                     # status; at most every 30 s
scontrol show job <id>                          # why pending / failed
scancel <id>                                    # cancel
\end{lstlisting}
